\chapter{Introduction}
\label{chap:introduction}

\section{Background to the Study}

\subsection{Historical Web Caching and Human Traffic Distributions}
Web caches are fundamental infrastructure systems designed to reduce web server load and latency by storing resources geographically and network-wise closer to the user \cite{breslau_web_1999}. By serving repeated requests from intermediate storage, caching prevents the need to fetch content directly from the origin web server on every individual request \cite{breslau_web_1999, umar_experimental_2026}. Historically, web caching algorithms have been highly effective because human browsing patterns naturally follow the Zipfian law, where a small set of popular web pages are requested frequently while the vast majority are accessed rarely \cite{breslau_web_1999}. Because human traffic exhibits high temporal locality—where recently accessed data is highly likely to be requested again in the near future—caching systems can easily store popular content at the network edge. This strong human locality of reference established the theoretical and practical foundation for traditional caching eviction policies, such as the Least Recently Used (LRU) algorithm \cite{breslau_web_1999}.

\subsection{The Proliferation of AI Scrapers and Crawlers}
However, the modern web ecosystem has experienced a massive influx of Artificial Intelligence (AI) scrapers and automated crawlers designed to harvest vast quantities of web data for training machine learning models and large language models (LLMs) \cite{kim_scrapers_2025, zhang_rethinking_2026}. While automated web crawlers have existed for decades—historically utilized for search engine indexing and search engine optimization (SEO)—the explosive growth of AI has fundamentally altered the threat landscape and the sheer scale of network traffic \cite{kim_scrapers_2025}. Today, aggressive AI data scrapers and live-retrieval crawlers overwhelm network edges to build massive training corpora or power retrieval-augmented generation (RAG) tasks \cite{zhang_rethinking_2026}.

\subsection{The Disruption of the Caching Ecosystem}
Unlike human users, these AI scrapers typically make high-volume, singular, and non-recurring requests across wide swaths of web pages that will never be requested again \cite{zhang_rethinking_2026}. This behavior fundamentally breaks the Zipfian distribution pattern that traditional caching algorithms rely upon. Automated bots generate what is known as a "scan" pattern: a request stream that touches a large number of distinct objects with little-to-no reuse in a short period, replacing the varied and deep browsing habits of humans with superficial single accesses \cite{zhang_rethinking_2026}. Consequently, current caching algorithms fail to handle this uniform, non-recurring traffic effectively. When traditional caches are subjected to these scan patterns, the algorithms are effectively flushed, causing cache hit ratios to decline steeply and leading to severe operational inefficiencies \cite{zhang_rethinking_2026}.

\subsection{The Evasion of Passive Defenses}
Furthermore, mitigating this aggressive automated traffic is severely complicated by the evasion tactics of modern bots. Recent empirical studies demonstrate that modern AI scrapers selectively ignore standard voluntary domain directives such as robots.txt, rendering passive exclusion policies highly unreliable \cite{kim_scrapers_2025}. When web administrators attempt to enforce active, friction-based defenses like CAPTCHAs or behavioral biometrics, modern AI systems quickly bypass them using sophisticated deep learning solvers and reinforcement learning agents designed to mimic human environmental characteristics \cite{ousat_broken_2026}. Because automated systems actively evade traditional security measures while simultaneously breaking the foundational assumptions of temporal locality, web caches remain critically vulnerable to AI-induced infrastructure degradation.

\section{Statement of the Problem}
The fundamental problem addressed by this study is that high-volume, non-recurring scan patterns generated by modern Artificial Intelligence (AI) scrapers induce severe cache pollution, forcing traditional eviction algorithms to prematurely discard frequently accessed human content which ultimately degrades system performance, inflates origin server loads, and increases latency for legitimate users \cite{zhang_rethinking_2026}.

\section{Aims and Objectives}

\subsection{Primary Aim}
The primary aim of this research is to design and implement a dynamic, heuristic caching architecture that mitigates Artificial Intelligence (AI) scraper interference and preserves web cache efficiency \cite{zhang_rethinking_2026}.

\subsection{Specific Objectives}
To achieve this aim, the specific objectives of the study are to:
\begin{itemize}
    \item Build a lightweight heuristic detector utilizing passive network analysis, specifically evaluating HTTP header anomalies and structural requests \cite{van_boxem_shy_2026}.
    \item Assign a real-time, continuous "AI-score" ranging from 0.0 to 1.0 to incoming HTTP requests in $\mathcal{O}(1)$ time complexity to probabilistically quantify automated bot origins \cite{chandra_composite_2026, van_boxem_shy_2026}.
    \item Develop a Dynamic Traffic-Aware Cache (DTAC) architecture that actively modifies standard Least Recently Used (LRU) eviction algorithms based on the assigned real-time AI-scores \cite{zhang_rethinking_2026}.
    \item Evaluate the DTAC prototype by benchmarking its cache hit rates directly against standard LRU performance under controlled, simulated workloads of mixed human and AI traffic \cite{breslau_web_1999, zhang_rethinking_2026}.
\end{itemize}

\section{Scope of the Research}

\subsection{Architectural Boundaries}
This study is strictly focused on server-side web caching infrastructure and routing optimization in the presence of automated scraping workloads \cite{umar_experimental_2026, zhang_rethinking_2026}. Architecturally, the research is bounded by the design and implementation of a two-stage system. This system comprises a Heuristic Proxy—featuring a State Manager and an AI-Score Generator—and a Dynamic Traffic-Aware Cache (DTAC) positioned upstream of an origin web server \cite{van_boxem_shy_2026, zhang_rethinking_2026}. The scope excludes the development or integration of client-side bot detection mechanisms, such as CAPTCHAs or behavioral biometrics, focusing entirely on passive, network-layer routing.

\subsection{Methodological Scope}
The methodological scope of this research is strictly constrained by the latency and concurrency requirements of edge-proxy environments, necessitating algorithms that guarantee $\mathcal{O}(1)$ execution time. Consequently, the proxy's threat evaluation is bounded to passive, Application Layer (OSI Layer 7) HTTP/HTTPS traffic analysis, utilizing Bayesian probability (Log-Odds classification) and bitwise Exponentially Weighted Moving Averages (EWMA) to assess behavioral metadata such as request velocity and structural anomalies. The integration of computationally heavy machine learning models, such as Deep Neural Networks (DNNs), directly on the critical path is explicitly excluded to avoid unacceptable latency spikes. Furthermore, the methodology is designed specifically to mitigate application-layer AI scrapers and cache-polluting botnets; thus, the mitigation of lower-level volumetric DDoS attacks (OSI Layers 3 and 4) falls outside the purview of this study, operating under the assumption that standard upstream infrastructure handles such threats.

To ensure system stability under internet-scale traffic, the methodology enforces strict bounded-memory constraints, intentionally abandoning unbounded state tracking and globally sorted cache queues. Client state tracking is limited to the probabilistic Space-Saving algorithm, while cache eviction relies entirely on lazy evaluation via randomized sampling (the ``Power of N Choices'') and lock-sharded metadata tracking utilizing a dual-stage Ghost Cache. The evaluative scope of this architecture is bounded to measuring cache hit-density, memory consumption, and read-path CPU overhead. Validation is conducted through simulated traffic models that combine Zipfian (heavy-tailed) human request distributions with high-velocity scraper traffic, assessing the system's ability to mathematically degrade the cache admission rate of automated swarms while preserving origin resources for legitimate users.

\section{Justification for the Research}

\subsection{The Failure of Current Mitigations}
The exponential growth of large language models and AI tools has driven a corresponding surge in web scraping activity \cite{zhang_rethinking_2026}. Because empirical evidence confirms that AI scrapers routinely ignore voluntary compliance mechanisms like robots.txt, web administrators can no longer rely on protocol-based exclusion to protect their infrastructure \cite{kim_scrapers_2025}. Furthermore, current active mitigation strategies, such as CAPTCHAs and behavioral biometrics, degrade user experience and are increasingly bypassed by modern automated solvers \cite{ousat_broken_2026}. While architectural solutions like physically partitioning hardware for different traffic types (split caches) offer isolation, they introduce severe resource fragmentation and economic inefficiencies that are impractical for most small-to-medium platforms \cite{zhang_rethinking_2026}.

\subsection{Protection of Infrastructure and Cost Efficiency}
This research is justified by its direct utility in safeguarding web performance \cite{umar_experimental_2026}. While network-level defenses such as hierarchical IP hashing exist to mitigate high-volume scraper strain \cite{hoetzlein_protecting_2025}, server-side caching systems themselves must be architecturally modernized to survive modern traffic anomalies. By introducing an AI-aware heuristic eviction mechanism, the proposed DTAC system proactively protects human traffic by preserving the temporal locality benefits of caching for actual end users \cite{breslau_web_1999, zhang_rethinking_2026}. Successful implementation will significantly improve overall infrastructure efficiency, achieving substantially higher cache hit rates for human traffic during periods of heavy scraper congestion, thereby saving critical server compute, bandwidth, and operational costs.